% ============================================================
%  ABSTRACT
% ============================================================

\begin{abstract}
Large language models (LLMs) have achieved remarkable performance across a broad range of natural
language processing tasks, yet their practical utility at the edge remains constrained by prohibitive
computational costs.
This paper presents \textbf{NanoChat 2.0}, a systematic study of whether
recent architectural and training innovations transfer meaningfully to a compact, GPT-2-scale
autoregressive language model.
Building on the NanoChat framework \parencite{nanochat}, we conduct controlled experiments with
four independent modifications: alternative activation functions (GeLU, GeLUSine, GELU-Sinc-Perturbation, GMTU, and Turbulent), multi-token prediction (MTP) loss, an Engram conditional memory
layer \parencite{cheng2026conditional}, and learned Attention Residuals \parencite{chen2026attnres}.
All experiments share the same FineWeb-Edu pretraining corpus and are evaluated using validation
bits-per-byte (BPB) and the CORE score across 22 diverse tasks.
Our key finding is that the NanoChat baseline with ReLU\(^2\) activation, already carefully
optimised, proves surprisingly hard to beat at a depth-12 scale of roughly 85 million parameters.
MTP offers a modest advantage on knowledge-intensive sub-tasks (ARC, MMLU) at the cost of
training throughput; Attention Residuals and Engram both underperform the baseline on aggregate
metrics.
Notably, enabling FP8 mixed-precision training yields the largest consistent gain across all
metrics with no architectural changes.
Our work serves as a reproducible proof-of-concept demonstrating that meaningful experimentation
on modern autoregressive LMs can be conducted on modest academic hardware, providing a controlled
setting to disentangle technique from scale.
\end{abstract}


% ============================================================
%  INTRODUCTION
% ============================================================
\section{Introduction}

\subsection{Motivation}

Large language Models have demonstrated strong performance across a wide range of natural language processing tasks, but training and deploying them requires substantial computational
resources, as modern frontier systems exceeds a hundred billions parameters
\parencite{deepseekai2025deepseekv3technicalreport}.
This makes it difficult to experiment with them in a small-scale setting that can be trained locally.
\textit{NanoChat} \parencite{nanochat} addresses this issue by providing a minimal, end-to-end chatbot framework designed to reproduce GPT-2-style language modelling behaviour at a much smaller scale and cost.
Its compact design makes it a useful experimental testbed: large enough to capture the essential components of modern autoregressive language models, yet small enough to be trained and modified under resources afforded in an academic setting. As such, NanoChat is ideal architecture for systematically evaluating whether recent advances in  model architecture and training translate to meaningful gains in practice.

A second area of impact is scientific reproducibility. As ablation studies on frontier models requires large compute, the academic research community has limited means to verify or improve on reported gains. A small, well-specified testbed like NanoChat allows controlled experiments in which a single
variable can be changed in isolation. This helps determine whether claimed improvements manifest from the techniques per se or emerge from scale.

\subsection{Objective}

Through this project we introduce \textbf{NanoChat 2.0}, an improved small-scale language model built on the NanoChat framework.
Our central research question is: \textit{to what extent can recent architectural and training innovations improve the efficiency and performance of a small-scale language model?}
We select four techniques that represent distinct axes of model improvement--activation function
design, training objective, memory augmentation, and residual connection design--and evaluate each modiciation in isolation against a carefully replicated baseline.
As a secondary goal, we study the behaviour of FP8 mixed-precision training as an orthogonal efficiency lever.

\subsection{NLP Task}

Large Language Models are commonly associated with broad downstream capabilities such as summarisation, machine translation, question answering, and sentiment analysis. These abilities are fundamentally grounded in the core task of autoregressive language modelling. Accordingly, our project focuses primarily on improving NanoChat as an \textbf{autoregressive language model for open-ended text generation}, using language modelling metrics and downstream
zero/few-shot task performance as proxies to measure general NLP capability.
